\section{A torsion diagnostic for finite-field recovery}\label{sec:torsion}

\begin{corollary}[Necessary torsion in bounded lifts]
\partformal{SmithRankLoss.lean}{smith_rank_loss,torsion_count_of_rank_and_nullity}{Existence of Smith normal form, the analytic rational-rank estimate and its asymptotic application, and the cokernel-torsion interpretation.}
\label{cor:dc-torsion}
Fix a prime $p$, integers $\Delta,t\ge1$, and $\varepsilon>0$.
Let $(E_n)$ satisfy the graph hypotheses of \cref{thm:dc-complex}.
Suppose matrices $M_n$ over $\mathbb F_p$ have diagonal blocks $I_t$,
off-diagonal support in $E_n$, and
$\dim\ker M_n\ge(1/2+\varepsilon)tN_n$.
For any integer lifts $\widetilde M_n$ with the same support restrictions,
literal diagonal blocks $I_t$, and uniformly bounded entries, at least
$(\varepsilon-o(1))tN_n$ of their \emph{nonzero} Smith invariant factors
are divisible by $p$.
\end{corollary}
\begin{informalproof}
Set $\rho_n=\operatorname{rank}_{\mathbb Q}\widetilde M_n$.
\cref{thm:dc-complex} gives $\rho_n\ge(1/2-o(1))tN_n$.
The unimodular transformations in Smith normal form remain invertible
modulo $p$.  Thus
\[
 \#\{\hbox{nonzero invariant factors divisible by }p\}
 =\rho_n-\operatorname{rank}_{\mathbb F_p}M_n
 \ge(\varepsilon-o(1))tN_n.
\]
Equivalently, the subgroup annihilated by $p$ in the integer cokernel
has at least this dimension over $\mathbb F_p$.  Zero invariant factors
give free abelian summands and are not counted as torsion.
\end{informalproof}

This conclusion permits such torsion; it is not a finite-field half-rate
theorem.  A fixed extension field can be represented over its prime field
by enlarging the fixed block width.  The analytic argument does not cover
growing block widths, unbounded complex coefficients, or nonlinear codes
on nonhyperfinite relations.  \cref{sec:mc-section} gives a separate
finite-field certificate on one existential expanding family, uniform
in the storage-code block width.  Neither result identifies all native
cone relations with that family or settles the unrestricted coding problem.
